\documentclass[12pt]{article}

%\usepackage[cm]{fullpage}
\setlength{\topmargin}{-1in}
\setlength{\oddsidemargin}{-0.25in}
\setlength{\evensidemargin}{-0.25in}
\setlength{\textwidth}{7in}
\setlength{\textheight}{10in}

\usepackage{array}
\usepackage{enumitem}
\usepackage{amssymb}
\usepackage{amsmath}
\usepackage{epstopdf}
\usepackage[T1]{fontenc}
\usepackage{graphics}
\usepackage{graphicx}
\usepackage{hyperref}
\usepackage{latexsym}
\usepackage{mathrsfs}
\usepackage{multicol}
\usepackage{setspace}
\usepackage{subfigure}
\usepackage{time}
\usepackage{times}

\newcommand{\To}{\Rightarrow}
\newcommand{\del}{\nabla}
\newcommand{\R}{\mathbb{R}}
\newcommand{\C}{\mathcal{C}}
\newcommand{\N}{\mathcal{N}}
\newcommand{\A}{\mathcal{A}}
\renewcommand{\O}{\mathcal{O}}
% \renewcommand{\N}{\mathbb{N}}
\newcommand{\ben}{\begin{enumerate}}
\newcommand{\een}{\end{enumerate}}
\newcommand{\eps}{\varepsilon}
\renewcommand{\epsilon}{\varepsilon}
\newcommand{\p}{\partial}
\newcommand{\<}{\langle}
\renewcommand{\>}{\rangle}
\renewcommand{\subset}{\subseteq}
\newcommand{\cont}{\Rightarrow \Leftarrow}
\newcommand{\back}{\backslash}
\newcommand{\norm}[1]{\|{#1}\|}
\newcommand{\abs}[1]{|{#1}|}
\newcommand{\ip}[1]{\langle{#1}\rangle}
\newcommand{\m}[1]{{\mathbf{#1}}}
\newcommand{\sts}{\m{S}^\textrm{T}\m{S}}
\newcommand{\rank}{\textrm{rank}}
\newcommand{\X}{\mathcal{X}}
\newcommand{\Y}{\mathcal{Y}}
\def\ds{\displaystyle}
\def\d{\partial}
\renewcommand{\arraystretch}{1.5}
\renewcommand{\L}{\mathscr{L}}

%\includegraphics[scale = 0.65, trim = 1in 3.25in 1in 3in]{blank_axes.pdf}


\pagestyle{empty}
\begin{document}
\noindent\rule{7in}{0.4pt}\\
\begin{center}
{\sc 
Differential Equations\\
Introduction to DEs	\\
}
\end{center}
\noindent\rule{7in}{0.4pt}
%%%%%%%%%%%%%%%%%%%%%%%%%%%%%%%%%%%%%%%%%%%%%%%%
\section*{What is a Differential Equation?}

\doublespacing
\noindent  Let us begin by defining what a differential equation (DE) is.
	\begin{quote}
	A {\bf differential equation} is an %equation 
	\line(1,0){200}\, that contains \\ \\ %a derivative or derivatives 
	\line(1,0){400}\, of \\ \\%some unknown function.
	\line(1,0){400}.
	\end{quote}
\singlespacing
\noindent  Let us look at a few examples to demonstrate what this means.\\

\noindent  Each of the following \line(1,0){200}: %{\bf is a DE}:
	\[ y'' + xy' = \sin x, \hspace{1cm}  \frac{dP}{dt} = P(P-10), \hspace{1cm} \ddot{Q} = 11Q. \]
\noindent  Note that there are three different notations for derivatives on display here.  Perhaps you
have not seen the dot notation in that third equation.  That does, in fact, mean differentiation, and
is sometimes used to indicate derivatives with respect to time.\\

\noindent  \line(1,0){500}:%\noindent  In contrast, none of the following is a DE:
	\[ \frac{d^2 y}{dx^2} + 2\frac{dy}{dx}, \hspace{1cm} x^2 + 3x = 19, \hspace{1cm} \sin y + e^x = 0\]
\noindent  The first is not an equation, so it certainly can't be a DE.  The other equations do not contain
derivatives, so they are not DEs.\\

\noindent   Let us look at one more example.

\end{document}